\section{Game-Theoretic Analysis}

\label{sec:game_theory}
% ══════════════════════════════════════════════════════════════════════════════

\subsection{Voter Strategic Behavior}

Under QV, strategic voting is bounded. Misrepresentation cost scales
quadratically:
\begin{equation}
  \text{Cost of misrepresentation} = \sum_j (k_{ij}^* - k_{ij})^2
\end{equation}

This makes strategic behavior increasingly expensive relative to sincere
voting as proposals increase.

\subsection{Collusion Resistance}

We implement pairwise penalization:
\begin{equation}
  V_j^{\text{adjusted}} = \sum_i k_{ij} \cdot \left(1 -
  \frac{1}{M} \sum_{i' \neq i} \text{sim}(i, i')^2\right)
\end{equation}

where $\text{sim}(i, i')$ measures voting pattern overlap.

\subsection{Reviewer Honesty}

\begin{proposition}[Honest Review Equilibrium]
Under the staking mechanism, honest reviewing is a Nash equilibrium when
expected quality from honest review exceeds that from strategic review
by more than the effort cost difference.
\end{proposition}

This holds because inter-reviewer consistency rewards agreement with
independent (likely honest) assessors, author feedback rewards genuine
engagement, and post-publication calibration penalizes divergence from
realized outcomes.


% ══════════════════════════════════════════════════════════════════════════════
